\subsection{Scott's cardinality in depth}
Scott's cardinality is a way to define the cardinality of a set that is independent from 
AC. We made use of the concept of rank in the definition of
cardinality. Specifically, the cardinality of a set $A$ is the set of all sets in
mutual injection with $A$ that also have minimal rank. Why is that? because if we 
were to consider the set of all sets in mutual injection with $A$ regardless of their rank, 
we would end up with an ill-formed set.

\begin{lemma}(Wrapping lemma)
    Let $A$ be a well-ordered infinite set, then, then $A \cup \{ X \}$ is 
    in bijection with $A$ for every possible choice of a generic set $X$.
\end{lemma}
\begin{proof}
    Let $\alpha$ be the ordinal in bijection with $A : \alpha \rightarrow (\alpha \cup \{X\})$ as follows:

    \begin{equation*}
        f(x) = 
        \begin{cases*}
            x' & if \( x = x' + 1 \)          \\
            \{X\}  & if \( x = 0 \)       \\
        \end{cases*}
    \end{equation*}                           \\

    We can exibit an inverse for this function herby proving its bijectivity by construction:

    \begin{equation*}
        f^{-1}(x) = 
        \begin{cases*}
            x + 1      & if \( x \in \alpha \)  \\
            0          & if \( x = \{X\} \) \\
        \end{cases*}
    \end{equation*}                             \\

    Since $A$ is an infinite well-ordered set, it is in bijection with some $\alpha$ hence some $(\alpha \;\cup \{X\})$ is in 
    bijection with $A$ as well. This result can be tought of as $|A| + 1 = |A|$ for every well ordered infinite set.
    

\end{proof}

\begin{lemma}(Ill-formedness of rankless definition of cardinality) 
    The class of every set in mutual injection with a given set is not a set for every 
    well-ordered infinite set.
\end{lemma}
\begin{proof}
    Let $A$ be a well-ordered infinite set, let $A^*$ be the class of all sets in mutual 
    injection with $A$. Notice that $A \cup \{X\}$ is in bijection with $A$ for every
    possible choice of a generic set $X$. Hence, it is in $A^*$. Now, consider the following construction:

    $$ A \cup \{X\} \in A^* $$ \vspace{0.05cm}
    $$ \{X\} \subseteq \bigcup_{S \in A^*} S $$
    $$ X \in P(\bigcup_{S \in A^*} S) $$
    
    This means that assuming that $A^*$ is a set, led us to a way of creating the universal set, 
    which itself gives us a way to fall within the russel paradox. This means that $A^*$ is not a 
    well-formed set, hence we can conclude that the class of all sets in mutual injection is not 
    a set for every well-ordered infinite set. Under AC, this generalizes to the fact that 
    no class of sets in bijection with any infinite sets can ever be a set. \\

\end{proof}

\newpage

One could ask himself why is that the case that the previous argument doesn't apply 
to the proper Scott's cardinality definition that takes the rank into account. The 
reason is simple: assuming that $S$ is a set of minimal rank, it's not the case that 
$ S \cup \{X\}$ is also of minimal rank. We can exibit a counterexample for this:

$$ V_{\omega} \subseteq V_{\omega} $$
$$ V_{\omega} \cup \{ V_{\omega} \} \not\subseteq V_{\omega} $$ \\

We took a set of rank $\omega$ of minimal rank $V_{\omega}$ and we 
constructed another set, $V_{\omega} \cup \{ V_{\omega} \}$, which is not of minimal rank while 
still being in bijection with $V_{\omega}$. This means that the previous argument will fail 
under the rank-aware definition of Scott's cardinality. \\

\begin{lemma}(Correctness of Scott's cardinality) 
    Even considering only sets of minimal rank, two sets in mutual injection
    still have the same cardinality regardless of their ranks.
\end{lemma}
\begin{proof}
    For every two sets $A$ and $B$, them being in mutual injection implies that they have the 
    same cardinality. In fact, Let $|A|$ be the set of all sets in mutual injection with 
    $A$ that also have minimal rank. May $|B|$ also be the set of all sets in mutual
    injection with $B$ that also have minimal rank. every set in $|A|$ is in mutual injection with 
    $A$, which is in mutal injection with $B$, which is in mutal injection with every set in $|B|$.

    $$ \forall X \in |A|\;\; X \darr A \land A \darr B \land B \darr Y \;\;\forall Y \in |B|$$
\end{proof}